\documentclass[10pt,a4paper]{amsart}
\usepackage{amsmath,amssymb,amsthm,graphicx}
% Retained minimal preamble for the Stacks Project chapter more-algebra.tex (More on Algebra).
% The chapter file pinned-original-source.tex carries only \input{preamble} on its first line, so no
% preamble text lives inside the chapter itself.  The Stacks Project uses one single shared preamble for
% every chapter; the snapshot stages it identically as sources/stacks-chow/preamble.tex and
% sources/stacks-etale-cohomology/preamble.tex.  The block below is copied mechanically, line for line,
% from sources/stacks-chow/preamble.tex lines 14-18 (the Xy-pic setup) and 154-217 (the native theorem
% declarations together with the mathematical macros).  Nothing was added, removed or redefined; the
% wrapper class and the amsmath/amssymb/amsthm/graphicx base packages are supplied by the assembler.
% Chapter: more-algebra.tex -- More on Algebra.
% For commutative diagrams we use Xy-pic
\usepackage[all]{xy}

% We use 2cell for 2-commutative diagrams.
\xyoption{2cell}

% Theorem environments.
%
\theoremstyle{plain}
\newtheorem{theorem}[subsection]{Theorem}
\newtheorem{proposition}[subsection]{Proposition}
\newtheorem{lemma}[subsection]{Lemma}

\theoremstyle{definition}
\newtheorem{definition}[subsection]{Definition}
\newtheorem{example}[subsection]{Example}
\newtheorem{exercise}[subsection]{Exercise}
\newtheorem{situation}[subsection]{Situation}

\theoremstyle{remark}
\newtheorem{remark}[subsection]{Remark}
\newtheorem{remarks}[subsection]{Remarks}

\numberwithin{equation}{subsection}

% Macros
%
\def\lim{\mathop{\mathrm{lim}}\nolimits}
\def\colim{\mathop{\mathrm{colim}}\nolimits}
\def\Spec{\mathop{\mathrm{Spec}}}
\def\Hom{\mathop{\mathrm{Hom}}\nolimits}
\def\Ext{\mathop{\mathrm{Ext}}\nolimits}
\def\SheafHom{\mathop{\mathcal{H}\!\mathit{om}}\nolimits}
\def\SheafExt{\mathop{\mathcal{E}\!\mathit{xt}}\nolimits}
\def\Sch{\mathit{Sch}}
\def\Mor{\mathop{\mathrm{Mor}}\nolimits}
\def\Ob{\mathop{\mathrm{Ob}}\nolimits}
\def\Sh{\mathop{\mathit{Sh}}\nolimits}
\def\NL{\mathop{N\!L}\nolimits}
\def\CH{\mathop{\mathrm{CH}}\nolimits}
\def\proetale{{pro\text{-}\acute{e}tale}}
\def\etale{{\acute{e}tale}}
\def\QCoh{\mathit{QCoh}}
\def\Ker{\mathop{\mathrm{Ker}}}
\def\Im{\mathop{\mathrm{Im}}}
\def\Coker{\mathop{\mathrm{Coker}}}
\def\Coim{\mathop{\mathrm{Coim}}}

% Boxtimes
%
\DeclareMathSymbol{\boxtimes}{\mathbin}{AMSa}{"02}

%
% Macros for moduli stacks/spaces
%
\def\QCohstack{\mathcal{QC}\!\mathit{oh}}
\def\Cohstack{\mathcal{C}\!\mathit{oh}}
\def\Spacesstack{\mathcal{S}\!\mathit{paces}}
\def\Quotfunctor{\mathrm{Quot}}
\def\Hilbfunctor{\mathrm{Hilb}}
\def\Curvesstack{\mathcal{C}\!\mathit{urves}}
\def\Polarizedstack{\mathcal{P}\!\mathit{olarized}}
\def\Complexesstack{\mathcal{C}\!\mathit{omplexes}}
% \Pic is the operator that assigns to X its picard group, usage \Pic(X)
% \Picardstack_{X/B} denotes the Picard stack of X over B
% \Picardfunctor_{X/B} denotes the Picard functor of X over B
\def\Pic{\mathop{\mathrm{Pic}}\nolimits}
\def\Picardstack{\mathcal{P}\!\mathit{ic}}
\def\Picardfunctor{\mathrm{Pic}}
\def\Deformationcategory{\mathcal{D}\!\mathit{ef}}
\usepackage{cleveref}
\crefname{originallemma}{original lemma}{original lemmas}
\Crefname{originallemma}{Original lemma}{Original lemmas}
\makeatletter
\newlabel{lemma-naive-derived-completion}{{091W}{source}{}{sourcepointer-1}{}}
\newlabel{lemma-naive-derived-completion@cref}{{[originallemma][0][]091W}{[1][0][]source}{}{}{}}
\makeatother
\crefname{originallemma}{original lemma}{original lemmas}
\Crefname{originallemma}{Original lemma}{Original lemmas}
\makeatletter
\newlabel{lemma-lift-complex-projectives}{{09AR}{source}{}{sourcepointer-2}{}}
\newlabel{lemma-lift-complex-projectives@cref}{{[originallemma][0][]09AR}{[1][0][]source}{}{}{}}
\makeatother
\crefname{originallemma}{original lemma}{original lemmas}
\Crefname{originallemma}{Original lemma}{Original lemmas}
\makeatletter
\newlabel{lemma-compute-Rlim-modules}{{091D}{source}{}{sourcepointer-3}{}}
\newlabel{lemma-compute-Rlim-modules@cref}{{[originallemma][0][]091D}{[1][0][]source}{}{}{}}
\makeatother
\crefname{originalremark}{original remark}{original remarks}
\Crefname{originalremark}{Original remark}{Original remarks}
\makeatletter
\newlabel{remark-how-unique}{{07KZ}{source}{}{sourcepointer-4}{}}
\newlabel{remark-how-unique@cref}{{[originalremark][0][]07KZ}{[1][0][]source}{}{}{}}
\makeatother
\crefname{originallemma}{original lemma}{original lemmas}
\Crefname{originallemma}{Original lemma}{Original lemmas}
\makeatletter
\newlabel{lemma-limit-flat}{{0912}{source}{}{sourcepointer-5}{}}
\newlabel{lemma-limit-flat@cref}{{[originallemma][0][]0912}{[1][0][]source}{}{}{}}
\makeatother
\title{A bounded above inverse system of derived objects with compatible reduction is derived complete and recovers each term by reduction}
\author{The Stacks Project Authors}
\date{stacks-project:more-algebra.tex:local-snapshot-20261012; GNU FDL 1.2 or later}
\begin{document}
\maketitle
\setcounter{subsection}{1}
\noindent\textit{Editorial note.} All delivered mathematical text is the author's own, copied from the pinned Stacks Project chapter file more-algebra.tex; no statement or proof step was rewritten, reordered or removed. The counted statement is the labeled lemma environment at source lines 28470-28483 (canonical tag 09AU) and the counted proof is the proof environment at source lines 28485-28517, which immediately follows it and is its own proof.\par\medskip
\noindent\textit{Editorial note.} The proof refers three times to labelled results of the same chapter (the criterion for derived completeness of a limit, the lifting of complexes, and the computation of a derived limit) and once jointly to a computation lemma and its uniqueness remark, and once more to the lemma on limits of flat modules. All are navigation pointers with their roles recorded in the external reference reasons.\par\medskip
\noindent\textit{Editorial note.} The sentence at source line 28468, recording that the author does not know whether the lemma holds for unbounded complexes, precedes the counted statement and lies outside the counted spans.\par\medskip
\noindent\textit{Editorial note.} The chapter file pinned-original-source.tex carries only an input of the shared preamble on its first line, so no preamble text lives inside the chapter. The retained file sources/stacks-more-algebra/retained-preamble.tex is a mechanical, line-for-line copy of the project's single shared preamble as staged in the snapshot (Xy-pic setup plus the native theorem declarations and mathematical macros). Nothing was added, removed or redefined, and the excerpt serves only as preamble, never as a span coordinate.\par\medskip
\noindent\textit{Editorial note.} The source URL is the source-level Stacks Project URL, as required when more than one card is produced from one source; the per-claim tag is carried by claim\_key and by this note. The counted statement and proof cite no bibliography entry, so this card needs no bibliography.\par\medskip
\noindent\textit{Allowed standard prerequisites (editorial).}\begin{itemize}
\item A derived limit of objects on which the elements of an ideal act nilpotently is derived complete, so the object K = R-lim K\_n is derived complete. used at proof lines 28486-28487. Gives the first conclusion of the lemma, namely that K is a derived complete object of the bounded below derived category.
\item A complex of modules over the quotient by a power of an ideal representing a given object can be chosen free with a prescribed vanishing bound, and matching maps between the reductions can be lifted to maps of complexes. used at proof lines 28490-28497. Produces the system of complexes P\_n of free modules with reductions isomorphic to the previous term.
\item The derived limit of the system is represented by the limit of the complexes, and the limit of a system of flat modules with surjective transition maps is flat. used at proof lines 28500-28511. Identifies K with the limit complex P and shows that P is a bounded above complex of flat modules, so that K tensor A over I\textasciicircum{}t is represented by P tensor A over I\textasciicircum{}t and the termwise limit can be computed.
\end{itemize}\medskip
\ifdefined\hypertarget\hypertarget{sourcepointer-1}{}\fi
\noindent\textit{Source reference.} Original lemma 091W: sources/stacks-more-algebra/pinned-original-source.tex, source lines 26288--26288. The reference supplies the criterion that a derived limit of objects killed by powers of the elements of I is derived complete, used at proof lines 28486-28487 for the first conclusion; the second conclusion, which is the substance of the proof, is obtained inside the counted span.\par\medskip
\ifdefined\hypertarget\hypertarget{sourcepointer-2}{}\fi
\noindent\textit{Source reference.} Original lemma 09AR: sources/stacks-more-algebra/pinned-original-source.tex, source lines 19469--19469. The reference supplies the lifting of maps of complexes along the reductions, used at proof lines 28492-28497 to build the system P\_n with strict reduction isomorphisms; the use made of that system is inside the counted span.\par\medskip
\ifdefined\hypertarget\hypertarget{sourcepointer-3}{}\fi
\noindent\textit{Source reference.} Original lemma 091D: sources/stacks-more-algebra/pinned-original-source.tex, source lines 23574--23574. The reference identifies the derived limit K with the limit complex P, invoked at proof lines 28500-28503; the subsequent flatness and stabilization computations are inside the counted span.\par\medskip
\ifdefined\hypertarget\hypertarget{sourcepointer-4}{}\fi
\noindent\textit{Source reference.} Original remark 07KZ: sources/stacks-more-algebra/pinned-original-source.tex, source lines 23753--23753. The remark is cited together with the previous lemma at proof lines 28502-28503 to justify representing the derived limit by the limit complex; it is a pointer accompanying that identification and carries no separate mathematical input.\par\medskip
\ifdefined\hypertarget\hypertarget{sourcepointer-5}{}\fi
\noindent\textit{Source reference.} Original lemma 0912: sources/stacks-more-algebra/pinned-original-source.tex, source lines 6475--6475. The reference is used twice, at proof lines 28506-28507 to conclude that the limit complex P is a bounded above complex of flat modules and at proof lines 28510-28511 to compute P tensor A over I\textasciicircum{}t termwise; both uses are inputs, and the stabilization argument built on them is inside the counted span.\par\medskip
\begin{lemma}
\label{lemma-Rlim-gives-complete}
Let $A$ be a ring and $I \subset A$ an ideal. Suppose
given $K_n \in D(A/I^n)$ and maps $K_{n + 1} \to K_n$ in
$D(A/I^{n + 1})$. If
\begin{enumerate}
\item $A$ is Noetherian,
\item $K_1$ is bounded above, and
\item the maps induce isomorphisms
$K_{n + 1} \otimes_{A/I^{n + 1}}^\mathbf{L} A/I^n \to K_n$,
\end{enumerate}
then $K = R\lim K_n$ is a derived complete object of $D^-(A)$ and
$K \otimes_A^\mathbf{L} A/I^n \to K_n$ is an isomorphism for all $n$.
\end{lemma}

\begin{proof}
The object $K$ of $D(A)$ is derived complete by
Lemma \ref{lemma-naive-derived-completion}.

\medskip\noindent
Suppose that $H^i(K_1) = 0$ for $i > b$. Then we can find a complex of
free $A/I$-modules $P_1^\bullet$ representing $K_1$ with $P_1^i = 0$
for $i > b$. By Lemma \ref{lemma-lift-complex-projectives}
we can, by induction on $n > 1$, find
complexes $P_n^\bullet$ of free $A/I^n$-modules representing $K_n$
and maps $P_n^\bullet \to P_{n - 1}^\bullet$ representing the maps
$K_n \to K_{n - 1}$ inducing isomorphisms (!)
of complexes $P_n^\bullet/I^{n - 1}P_n^\bullet \to P_{n - 1}^\bullet$.

\medskip\noindent
Thus we have arrived at the situation where $R\lim K_n$ is represented by
$P^\bullet = \lim P_n^\bullet$, see
Lemma \ref{lemma-compute-Rlim-modules} and
Remark \ref{remark-how-unique}.
The complexes $P_n^\bullet$ are uniformly bounded above complexes
of flat $A/I^n$-modules and the transition maps are termwise surjective.
Then $P^\bullet$ is a bounded above complex of flat $A$-modules by
Lemma \ref{lemma-limit-flat}.
It follows that $K \otimes_A^\mathbf{L} A/I^t$ is represented by
$P^\bullet \otimes_A A/I^t$. We have
$P^\bullet \otimes_A A/I^t = \lim P_n^\bullet \otimes_A A/I^t$
termwise by Lemma \ref{lemma-limit-flat}. The transition maps
$P_{n + 1}^\bullet \otimes_A A/I^t \to P_n^\bullet \otimes_A A/I^t$
are isomorphisms for $n \geq t$ by our choice of $P_n^\bullet$, hence we have
$\lim P_n^\bullet \otimes_A A/I^t = P_t^\bullet \otimes_A A/I^t = P_t^\bullet$.
Since $P_t^\bullet$ represents $K_t$, we see that
$K \otimes_A^\mathbf{L} A/I^t \to K_t$ is an isomorphism.
\end{proof}

\end{document}
